Um die Geschwindigkeit der Elektronen und Positronen vor der Kollision zu bestimmen, verwenden wir die Energie-Impuls-Beziehung in der relativistischen Physik. Die Gesamtenergie im Schwerpunktsystem ist gegeben durch die Ruhemasse des Z-Bosons, also \( E_{\text{cm}} = m_Z c^2 = 91 \, \text{GeV} \).

Da die Elektronen und Positronen gleich beschleunigt werden, haben sie im Schwerpunktsystem gleiche Geschwindigkeiten, aber entgegengesetzte Impulse. Die Gesamtenergie eines Teilchens ist gegeben durch

\[
E = \gamma m_e c^2
\]

wobei \(\gamma\) der Lorentzfaktor ist, definiert als \(\gamma = \frac{1}{\sqrt{1 - \frac{v^2}{c^2}}}\), und \(m_e\) die Ruhemasse des Elektrons (bzw. Positrons) ist. Da die Gesamtenergie im Schwerpunktsystem gleich der Summe der Energien der beiden Teilchen ist, haben wir

\[
2 \gamma m_e c^2 = m_Z c^2
\]

Daraus folgt

\[
\gamma = \frac{m_Z}{2 m_e}
\]

Die Geschwindigkeit \(v\) der Teilchen ergibt sich dann aus der Beziehung für den Lorentzfaktor:

\[
v = c \sqrt{1 - \frac{1}{\gamma^2}}
\]

Für den Fall, dass ein Elektron auf ein ruhendes Positron beschleunigt wird, muss die Gesamtenergie im Schwerpunktsystem ebenfalls \(m_Z c^2\) betragen. Die Energie des ruhenden Positrons ist einfach \(m_e c^2\), während die Energie des bewegten Elektrons \(E = \gamma m_e c^2\) ist. Die Bedingung für die Gesamtenergie ist dann

\[
\gamma m_e c^2 + m_e c^2 = m_Z c^2
\]

was zu

\[
\gamma = \frac{m_Z + m_e}{m_e}
\]

führt. Die Geschwindigkeit des Elektrons ist dann ebenfalls gegeben durch

\[
v = c \sqrt{1 - \frac{1}{\gamma^2}}
\]

Diese Gleichungen erlauben es, die Geschwindigkeit der Teilchen in beiden Szenarien zu berechnen.